% A.64 reserved by the root.
\subsection{A normalized prime-power coordinate at the dyadic boundary}

\begin{statement}{A.64}{Proved}
{The parameters $\pm16p^n$ are recognized by their full congruence data}
Let $p$ be an odd prime, $n\ge1$, and
$\lambda_0=\epsilon16p^n$ with $\epsilon\in\{1,-1\}$.
Every noncuspidal Selmer section whose full integral Legendre
representation is conjugate to that of $s_{\lambda_0}$ by an
element of $H$ is precisely $s_{\lambda_0}$.
The conclusion also holds for the six fractional-linear
transforms of $\lambda_0$ permuting $0,1,\infty$.
The geometric-conjugator hypothesis is retained here.
\end{statement}
\begin{proof}
The realizing Tate system gives the finite-ramification
conditions A.14.  Let $A,B$ be the two completed Kummer
coordinates.  The map to $H^{\rm ab}$ in A.56 contains
the separate mod-eight characters of $A$ and $B$.
Conjugate arithmetic sections have the same classes in this
abelian extension.  Therefore their valuation classes at two
satisfy
\[
                (v_2(A),v_2(B))\equiv(4,0)\pmod8.
\tag{A.64.1}
\]
These are congruences in $\widehat\Z/8$ obtained from the
actual common quotient, not from a choice of an arbitrary
linear intertwiner.

Since $\lambda_0$ is divisible by $16$ with exact
valuation four, $\lambda_0^2-\lambda_0+1$ and
$1-\lambda_0$ are odd.  The Legendre formula gives
\[
                         v_2(j(E_{\lambda_0}))=8-2\cdot4=0.
\tag{A.64.2}
\]
Thus the realized curve is potentially good at two, and
A.37 makes $s_2$ an ordinary point $x_2$.
Its Legendre curve is also potentially good, by the local
Tate-system comparison and potential good reduction.
The valuation pairs of an ordinary point of $U(\Q_2)$
have the form
\[
       (k,0),\quad(0,k),\quad(-k,-k),\qquad k>0.
\tag{A.64.3}
\]
Indeed the valuations of $x_2$ and $1-x_2$ obey these
alternatives, and both cannot be zero because both residues
would be one in $\mathbf F_2$.
In each alternative the Legendre formula gives
$v_2(j)=8-2k$.  Potential good reduction therefore forces
$1\le k\le4$.
Among these twelve integer pairs, only $(4,0)$ has
residue $(4,0)$ modulo eight: the pair $(-4,-4)$ has
residue $(4,4)$, and the other positive-coordinate and
negative-coordinate alternatives have different residues.
Consequently
\[
                          v_2(A)=4,\qquad v_2(B)=0
                              \quad\hbox{in }\widehat\Z.
\tag{A.64.4}
\]

At an odd prime $r\ne p$, the parameter $\lambda_0$ is
a unit.  If $1-\lambda_0$ is a unit, the local point is
integral; otherwise the exact labeled cyclotomic-line
comparison A.41 gives contact label one, for either an
ordinary point or a cusp section.  Hence $v_r(A)=0$
at every such $r$.
At $p$ the same comparison gives contact label zero,
since $v_p(\lambda_0)=n>0$.
Combining these facts with (A.64.4) in the sign-and-valuation
description of the completed multiplicative group gives
\[
                         A=\epsilon'16p^\beta,
            \qquad \epsilon'\in\{1,-1\},\quad
                         \beta\in\widehat\Z.
\tag{A.64.5}
\]

Suppose $s_p$ were cuspidal.  Its label is zero, and the
exact formal Tate class in A.63 is
\[
                         A^2/256=p^{2\beta}.
\tag{A.64.6}
\]
At coefficient $p$ this lies on the rational Kummer line
$\Q_p\kappa_p(p)$, which A.63.4 excludes for the Tate
class of the realized elliptic curve, even after a common
finite splitting extension and with an arbitrary scalar
intertwiner.  Thus $s_p$ is an ordinary point.

Its contact valuation is an ordinary positive integer $k$.
The full local valuation map on (A.64.5) gives
$\beta=k$ in $\widehat\Z$, so $A=\epsilon'16p^k$
is one rational number.  At the density-one set of ordinary
point localizations supplied by A.14, injectivity of
$\Q_r^\times\to\widehat{\Q_r^\times}$ makes every
local parameter equal to this rational number.
The finite-support recognition argument used in A.14 and
A.37 therefore makes $s$ its point section.
Finally, the primitive-isogeny and ordered-two-torsion
argument at the end of A.54 identifies that point's
parameter exactly with $\lambda_0$.  This step proves the
claimed parameter equality without presupposing an equality
of contact valuations before the point is recognized.

The fractional-linear extension uses the actual substitutions
in A.63.  To check preservation of geometric conjugacy as
well, choose a geometric element $\gamma\in\Pi$ whose
Legendre image is the given conjugator in $H$; this is
possible since $H$ is the compact image of $\Pi$.
For $t\mapsto1-t$, the constant quadratic twist and its
change of basis take that image to a geometric conjugator
for the transformed family.  For $t\mapsto1/t$, the image
of $\gamma$ gains the scalar factor $\chi_t(\gamma)$
from (A.63.8), in addition to the same change of basis;
here $\chi_t$ is the quadratic character of the cover $z^2=t$.
This factor is $1$ or $-1$ and is central, so it does not
change the conjugacy equality.
On the two arithmetic sections the twisting characters
agree by A.16, as proved in A.63.
Thus the conjugator is the actual geometric image of
the transformed $\gamma$, and belongs to the transformed
geometric closure.  Applying the result and then the
inverse automorphism proves the permutation assertion.
\end{proof}
